Accurate material forecasting is a critical operational challenge in apparel manufacturing. Discrepancies between planned Bill of Materials (BOM) estimates and actual fabric consumption routinely generate procurement errors of 8--12\%, leading to overstock, emergency re-orders, and supply chain inefficiencies that inflate both cost and environmental impact. This thesis designs, evaluates, and deploys a comprehensive machine learning-based fabric consumption forecasting system as a data-driven replacement for the static BOM formula.

A physics-informed synthetic dataset generator produces 5{,}000 production orders spanning January~2022 to September~2024, covering five garment types (T-Shirt, Shirt, Pants, Dress, Jacket), five fabric categories, four production seasons, and nine operationally observable input features. Four models are trained under identical preprocessing and evaluation conditions on a strict 64/16/20 train-validation-test split: Linear Regression, Random Forest, XGBoost, and Long Short-Term Memory (LSTM). A dynamic inverse-RMSE\textsuperscript{2} weighted-average ensemble combines their predictions.

On the held-out 1{,}000-order test set, the LSTM achieves RMSE 346\,m (379\,yd), MAPE 5.14\%, and R\textsuperscript{2}=0.9955---the best single model, representing a 21.8\% RMSE reduction over the BOM physics baseline (484\,yd). XGBoost (RMSE 574\,yd, MAPE 8.14\%, R\textsuperscript{2}=0.9897) provides the strongest tree-based performance. The ensemble (RMSE 422\,yd, R\textsuperscript{2}=0.9944, MAPE 5.92\%) offers robust procurement-grade prediction across the order-size distribution. Variance analysis reveals a size-dependent directional bias: small orders over-consume by $+$4.49\% while large orders under-consume by $-$2.79\% relative to the BOM plan. Feature importance analysis identifies garment type (0.518) and order quantity (0.297) as jointly accounting for 81.5\% of XGBoost's predictive gain. Winter is the most challenging season (MAPE 9.67\%) and polyester the most challenging fabric category (MAPE 10.31\%). Five-fold cross-validation confirms that test-set performance is stable and generalisable.

At a reference volume of 1{,}000 orders per year, the LSTM's 21.8\% RMSE improvement translates to \$180{,}000 in annual waste cost savings, a payback period of under four months on a \$50{,}000 implementation investment, and a five-year NPV of \$730{,}000 at a 5\% cost of capital. Environmental savings amount to approximately 6{,}850\,kg of avoided CO\textsubscript{2} emissions, 103{,}500 litres of water conserved, and 4{,}300\,kg of fabric diverted from landfill. All trained models are deployed in a six-module Streamlit web application providing single-order prediction, batch processing, 90\% confidence intervals, and an integrated ROI calculator.

\medskip
\begin{sloppypar}
\noindent\textbf{Keywords:} Material Forecasting, Apparel Manufacturing, Machine Learning, LSTM, XGBoost, BOM, Fabric Consumption, Ensemble Learning, Streamlit
\end{sloppypar}
